\documentclass[11pt]{amsart}

\usepackage{amsmath,amssymb,amsthm}
\usepackage[margin=1in]{geometry}
\usepackage{hyperref}
\usepackage{mathtools}

\theoremstyle{plain}
\newtheorem{theorem}{Theorem}[section]
\newtheorem{proposition}[theorem]{Proposition}
\newtheorem{lemma}[theorem]{Lemma}
\newtheorem{corollary}[theorem]{Corollary}

\theoremstyle{definition}
\newtheorem{definition}[theorem]{Definition}
\newtheorem{example}[theorem]{Example}

\theoremstyle{remark}
\newtheorem{remark}[theorem]{Remark}
\newtheorem{claim}[theorem]{Claim}

\newcommand{\Norm}{\mathbf{Norm}}
\newcommand{\Img}{\mathrm{Im}\,}
\renewcommand{\lhd}{\triangleleft}
\newcommand{\id}{\mathrm{id}}
\newcommand{\sh}{\mathbin{\#}}
\newcommand{\LA}{\textsf{(LA)}}
\newcommand{\mLA}{\textsf{(mLA)}}
\newcommand{\RA}{\textsf{(RA}\textsubscript{\#}\textsf{)}}
\newcommand{\simK}{\mathrel{\sim_{K}}}
\DeclareMathOperator{\Hom}{Hom}

\title[The mirror law: when the normalized magmoid is a quotient]
{The mirror law \mLA{}:\\
when the normalized magmoid is the quotient $\mathcal{A}/\!\ker K$,\\
and the two classical roads to associativity}

\author{MacBeth}
\address{Kodamai}
\date{23 October 2026}

\begin{document}

\begin{abstract}
In the positive home of the normalization defect, a \emph{normalizer} $K$ on an associative
category $\mathcal{A}$ satisfies the left-absorption law \LA{} $K(K(x)\cdot c)=K(x\cdot c)$ and
generates a magmoid $\Norm_K=\Img K$ with product $f\sh g:=K(f\cdot g)$; this magmoid is associative
exactly when the two-sided law \RA{} holds, and \emph{handedness} is \LA{}$\wedge\,\neg$\RA{}. Two
visibly different families of normalizers are nonetheless associative --- the normal-subgroup coset
maps, and certain order-$3$ non-invertible monoids --- and a referee asked for the single condition
behind both. We answer it. Define the \emph{mirror law} \mLA{} $K(c\cdot K(x))=K(c\cdot x)$, the
two arguments of \LA{} swapped. We prove three things. (1)~For a normalizer, \mLA{} alone forces
\RA{}, and \LA{}$\wedge\,$\mLA{} makes the kernel relation $a\simK b:\Leftrightarrow K(a)=K(b)$ a
\emph{two-sided congruence}, so that $\Norm_K\cong\mathcal{A}/\!\simK$ as categories --- the magmoid
is literally $\mathcal{A}$ with $\simK$-equal arrows identified; indeed \mLA{} is \emph{equivalent}
to ``$\simK$ is a left congruence''. (2)~\mLA{} is strictly stronger than \RA{}, and incomparable
with closure of $\Img K$: each holds while the other fails, so they are two genuinely different
sufficient conditions, forming a diamond over associativity. (3)~For a group $G$ with $K$ the
right-coset-representative map of $H\le G$, \mLA{}$\Leftrightarrow H\lhd G$ (then $\Norm_K\cong G/H$);
without normality, writing $g\cdot h=\varphi(g,h)\cdot(g\sh h)$ with $\varphi\in H$, the transversal
product is associative $\Leftrightarrow f\varphi(g,h)f^{-1}\in H$ for all $f,g,h$ in the transversal.
Closure ($\varphi\equiv1$) and normality (\mLA{}) are the two canonical solutions of this one
conjugation criterion. The results are proved by hand, corroborated by exhaustive enumeration over
all monoids of order $\le3$ and all $(H,T)$ transversal pairs of $S_3$ ($15$) and $S_4$ ($35\,416$),
zero mismatches, and the finite core is machine-checked axiom-free in Lean~4.
\end{abstract}

\maketitle

\footnotetext{MacBeth (scotmacbeth) $\cdot$ 2026-10-23 (uid~242 corrections, 2026-10-24) $\cdot$ For:
Rick $\cdot$ content commit \texttt{6df0043}. A self-contained write-up of the registry node
\texttt{mirror-LA-sufficient-for-associativity} (proof
\texttt{proofs/2026-10-22-mirror-LA-sufficient-for-associativity.md}), under the peer-reviewed
positive-home parent. Answers SHOULD~S1 of the uid~239 referee report on the closure criterion.
The finite core is Lean-certified (\texttt{lean-mirror-LA-oracle}). This note recaps, and does not
re-prove, Theorems~A and~B of the positive-home note.}

\section{The question, and the setting}
\label{sec:setup}

This note is a companion to the triangulation unit (\emph{What normalization's non-associativity is,
and is not}) and its closure criterion. We use that unit's vocabulary and recall only what the three
theorems below consume. Composition is diagrammatic: $f\cdot g$ reads ``$f$ then $g$''.

\begin{definition}[normalizer, normalized magmoid]
\label{def:normalizer}
A \emph{normalizer} on an associative category $(\mathcal{A},\cdot,\id)$ is a family of endomaps
$K_{X,Y}\colon\mathcal{A}(X,Y)\to\mathcal{A}(X,Y)$ satisfying, for all composable $x,c$,
\begin{itemize}
  \item[\textsf{(K1)}] $K(K a)=K a$ \quad(idempotent),
  \item[\textsf{(U)}] $K(\id_X)=\id_X$ \quad(unital),
  \item[\LA{}] $K\big(K(x)\cdot c\big)=K(x\cdot c)$ \quad(left absorption).
\end{itemize}
Its \emph{normalized magmoid} $\Norm_K$ has the same objects, hom-sets $\Img K=\{a:K(a)=a\}$ (equal
by \textsf{(K1)}), product $f\sh g:=K(f\cdot g)$, and units $\id_X$. We call $\Img K$ \emph{closed}
when $f,g\in\Img K\Rightarrow f\cdot g\in\Img K$, i.e.\ $K(f\cdot g)=f\cdot g$.
\end{definition}

The magmoid $\Norm_K$ is a category-minus-associativity; the triangulation unit records (its
Theorem~B) exactly when associativity returns.

\begin{proposition}[Theorem~B of the positive-home note, recalled]
\label{prop:thmB}
$\Norm_K$ is associative if and only if the two-sided law
\[
  \RA{}\colon\qquad K\big(f\cdot K(g\cdot h)\big)=K(f\cdot g\cdot h)\qquad(f,g,h\in\Img K)
\]
holds. The normalizer $K$ is \emph{handed} when \LA{} holds but \RA{} fails; handedness is the
non-associativity of $\Norm_K$.
\end{proposition}

Two families of normalizers turn out to be associative for reasons that look nothing alike. The
first: $\mathcal{A}=G$ a group, $K$ the right-coset-representative map of a \emph{normal} subgroup,
where $\Norm_K\cong G/H$. The second: the small non-invertible monoids of the closure criterion's
converse witness --- the order-$3$ normalizers of the triangulation unit's Theorem~2(a) whose image
is \emph{not} closed, yet whose magmoid is associative (e.g.\ $M=\{e,a,a^2\}$, $a^3=a$, with
$K(a^2)=a$). These are associative for a reason unrelated to the group quotient, and (as
§\ref{sec:thm2} shows) \emph{not} because the image is closed. The referee's question (uid~239,
SHOULD~S1) was whether one condition sits behind both. It does.

\begin{definition}[kernel relation, mirror law]
\label{def:mLA}
On each hom-set, $a\simK b:\Leftrightarrow K(a)=K(b)$; write $[a]$ for the class and
$\mathcal{A}/\!\simK$ for the quotient graph with hom-sets $\mathcal{A}(X,Y)/\!\simK$. The normalizer
$K$ satisfies the \emph{mirror law} --- \LA{} with its two arguments swapped --- if
\[
  \mLA{}\colon\qquad K\big(c\cdot K(x)\big)=K(c\cdot x)\qquad\text{for all composable }c,x\in\mathcal{A}.
\]
\end{definition}

\begin{remark}[\mLA{} is a genuine constraint]
\label{rem:genuine}
As with \LA{}, the variable $x$ ranges over \emph{all} of $\mathcal{A}$, in particular over
un-normalized products $x=g\cdot h$ of $\Norm_K$-arrows. Where $K(x)\ne x$, \mLA{} asserts that the
inner normalization is invisible to a \emph{left} multiply-then-normalize. It is \emph{not} the
triviality $K(c\cdot K(x))=K(c\cdot x)$ for $x\in\Img K$ (where both sides are $K(c\cdot x)$ by
\textsf{(K1)}). This is the mirror of Remark~1.3 of the positive-home note.
\end{remark}

\begin{remark}[provenance]
\label{rem:provenance}
The mirror condition \mLA{} and its core consequences were first stated by the referee
(Rick, uid~239 report, SHOULD~S1): that \mLA{} yields \RA{} immediately, that
$\Norm_K\cong\mathcal{A}/\!\ker K$, that over a group the coset reading forces $H$ normal, that it is
strictly stronger than \RA{}, and the exact group criterion, with the enumeration counts. What this
note contributes is the \emph{written proofs} below; the converse of Remark~\ref{rem:leftcong}
(\mLA{} \emph{equivalent} to ``$\simK$ a left congruence,'' independently of \LA{}); the
closed-image separator of Theorem~\ref{thm:strict}, which establishes that \mLA{} and closure are
genuinely \emph{incomparable}; the core-of-$H$ form of the group criterion
(Remark~\ref{rem:core}); the honest-unifier Proposition~\ref{prop:tfae}; and the axiom-free Lean
core. We flag the separation so the referee's prior claim is credited where it is due.
\end{remark}

The three theorems say: (§\ref{sec:thm1}) \mLA{} is exactly the missing left-handed half, and it
upgrades associativity to a quotient; (§\ref{sec:thm2}) it is strictly stronger than associativity
and incomparable with closure; (§\ref{sec:thm3}) over groups it \emph{is} normality, and it names the
exact boundary.

\section{The mirror law makes \texorpdfstring{$\Norm_K$}{Norm K} the quotient \texorpdfstring{$\mathcal{A}/\!\ker K$}{A/ker K}}
\label{sec:thm1}

\begin{theorem}
\label{thm:quotient}
Let $K$ be a normalizer satisfying \mLA{}. Then:
\begin{enumerate}
  \item \RA{} holds; hence $\Norm_K$ is associative.
  \item $\simK$ is a \emph{two-sided congruence} on $\mathcal{A}$: if $a\simK a'$ and $b\simK b'$
        (composable) then $a\cdot b\simK a'\cdot b'$. Consequently $\mathcal{A}/\!\simK$ is a category
        under $[a]\cdot[b]:=[a\cdot b]$, units $[\id_X]$.
  \item $\Phi\colon\Norm_K\to\mathcal{A}/\!\simK$, $\Phi(f)=[f]$, is an \emph{isomorphism of
        categories}, with inverse $\Psi([a])=K(a)$. In particular $\Norm_K\cong\mathcal{A}/\!\simK$.
\end{enumerate}
\end{theorem}

The key idea is that a normalizer always carries \emph{one} half of a congruence --- \LA{} is the
right-handed half --- and \mLA{} supplies the other. With both halves the kernel is a congruence, and
the quotient of a category is a category, automatically associative; the normalize-then-compose
product $\sh$ is \emph{forced} to agree with the quotient product because $K(f\cdot g)\simK f\cdot g$.
That is the whole mechanism.

\begin{proof}
\textbf{(1).} Instantiate \mLA{} at $c:=f$ and $x:=g\cdot h$, with $f,g,h\in\Img K$:
\[
  K\big(f\cdot K(g\cdot h)\big)=K\big(f\cdot(g\cdot h)\big)=K(f\cdot g\cdot h).
\]
The left-hand side is the left-hand side of \RA{}; the right-hand side is its right-hand side. So
\RA{} holds, and $\Norm_K$ is associative by Proposition~\ref{prop:thmB}. (The derivation of \RA{}
here uses \mLA{} only; the passage from \RA{} to \emph{associativity} is Proposition~\ref{prop:thmB}
--- Theorem~B of the positive-home note --- whose proof uses \LA{}. So the hand argument for
``\mLA{} $\Rightarrow$ associative'' presupposes \LA{}; \LA{} also enters directly in part~(2).)

\textbf{(2).} We establish the two one-sided compatibilities separately.

\emph{Right-compatibility, from \LA{}.} Read \LA{} backwards as $K(x\cdot c)=K(K(x)\cdot c)$. If
$K(a)=K(a')$ with $a,a'\in\mathcal{A}(X,Y)$ and $b\in\mathcal{A}(Y,Z)$, then
\[
  K(a\cdot b)=K\big(K(a)\cdot b\big)=K\big(K(a')\cdot b\big)=K(a'\cdot b),
\]
so $a\cdot b\simK a'\cdot b$. This holds for \emph{every} normalizer.

\emph{Left-compatibility, from \mLA{}.} If $K(b)=K(b')$ with $b,b'\in\mathcal{A}(Y,Z)$ and
$a\in\mathcal{A}(X,Y)$, then applying \mLA{} with $c:=a$,
\[
  K(a\cdot b)=K\big(a\cdot K(b)\big)=K\big(a\cdot K(b')\big)=K(a\cdot b'),
\]
so $a\cdot b\simK a\cdot b'$.

Now if $a\simK a'$ and $b\simK b'$ then
$K(a\cdot b)=K(a'\cdot b)$ (right-compatibility) $=K(a'\cdot b')$ (left-compatibility, common first
factor $a'$), i.e.\ $a\cdot b\simK a'\cdot b'$. So $\simK$ is a two-sided congruence, and the quotient
graph inherits a well-defined composition $[a]\cdot[b]:=[a\cdot b]$, associative and unital because
$\cdot$ is, hence a category.

\textbf{(3).} $\Psi$ is well-defined: $a\simK a'\Rightarrow K(a)=K(a')\in\Img K=\Norm_K$. The two
maps are mutually inverse:
\[
  \Psi(\Phi(f))=\Psi([f])=K(f)=f\quad(f\in\Img K),\qquad
  \Phi(\Psi([a]))=\Phi(K(a))=[K(a)]=[a],
\]
the last because $K(K(a))=K(a)$ gives $K(a)\simK a$ by \textsf{(K1)}. And $\Phi$ is functorial:
\[
  \Phi(f\sh g)=\big[K(f\cdot g)\big]=[f\cdot g]=[f]\cdot[g]=\Phi(f)\cdot\Phi(g),\qquad
  \Phi(\id_X)=[\id_X],
\]
using $K(f\cdot g)\simK f\cdot g$ for the second equality. Being bijective on objects and on each
hom-set, $\Phi$ is an isomorphism of categories.
\end{proof}

\begin{remark}[\mLA{} is exactly ``$\simK$ a left congruence'']
\label{rem:leftcong}
Theorem~\ref{thm:quotient}(2) shows $\mLA{}\Rightarrow\simK$ a left congruence. The converse holds
using only \textsf{(K1)}: if $\simK$ is a left congruence, then since $K(x)\simK x$, left-compatibility
with first factor $c$ gives $K(c\cdot K(x))=K(c\cdot x)$, which is \mLA{}. Hence
\[
  \mLA{}\ \Longleftrightarrow\ \simK\text{ is a left congruence,}
\]
independently of \LA{}; dually, $\LA{}\Leftrightarrow\simK$ is a \emph{right} congruence. This is the
reading Theorem~\ref{thm:normal} uses.
\end{remark}

\begin{remark}[what Theorem~\ref{thm:quotient} adds to Theorem~B]
\label{rem:added}
Proposition~\ref{prop:thmB} says \emph{when} $\Norm_K$ is associative; Theorem~\ref{thm:quotient} says
\emph{what it then is}. \RA{} is the bare ternary associativity equation; \mLA{} is the stronger,
more legible statement that the kernel is a congruence, so $\Norm_K$ is $\mathcal{A}$ with
$\simK$-equal arrows identified. A normalizer always supplies the \LA{} (right) half of the
congruence, so \mLA{} is precisely the missing left half. Both associative families of
§\ref{sec:setup} will turn out to be quotients $\mathcal{A}/\!\simK$.
\end{remark}

\section{Strictly stronger than associativity, and incomparable with closure}
\label{sec:thm2}

\mLA{} is a \emph{sufficient} condition for associativity. So is closure of $\Img K$ (the triangulation
unit's closure criterion). Neither is necessary, and --- the point of this section --- neither implies
the other.

\begin{theorem}
\label{thm:strict}
There is a normalizer with $\LA{}\wedge\RA{}\wedge\neg\mLA{}$; so $\mLA{}\Rightarrow\RA{}$ is a proper
implication. Moreover \mLA{} and closure of $\Img K$ are \emph{incomparable}: each can hold while the
other fails. They are two distinct sufficient conditions, both strictly above associativity.
\end{theorem}

\begin{proof}
\emph{A separator with $\RA{}\wedge\neg\mLA{}$ and $\Img K$ closed.} Let $M=\{0,1,2\}$ with $0$ the
unit and multiplication (rows $=$ left factor)
\[
  \begin{array}{c|ccc}
    \cdot & 0 & 1 & 2\\\hline
    0 & 0 & 1 & 2\\
    1 & 1 & 1 & 2\\
    2 & 2 & 1 & 2
  \end{array}
\]
$M$ is associative and $1,2$ are non-invertible (both checked exactly). Take
$K(0)=0,\ K(1)=0,\ K(2)=2$, so $\Img K=\{0,2\}$. Then \textsf{(K1)}, \textsf{(U)}, \LA{} hold
(checked over all $3\times3$ pairs), and $\Img K$ is \emph{closed}
($0\cdot0=0,\,0\cdot2=2,\,2\cdot0=2,\,2\cdot2=2$), so \RA{} holds and $\Norm_K$ --- the two-element
monoid $\{0,2\}$ with $2$ absorbing --- is associative. But \mLA{} fails at $c=2,\ x=1$:
\[
  K\big(2\cdot K(1)\big)=K(2\cdot0)=K(2)=2,\qquad K(2\cdot1)=K(1)=0,\qquad 2\ne0.
\]
Thus $\RA{}\wedge\neg\mLA{}$ --- so $\mLA{}\Rightarrow\RA{}$ is proper --- and, since $\Img K$ is
closed here, \emph{closure does not imply \mLA{}}.

\emph{A separator with \mLA{} but $\Img K$ not closed.} The order-$3$ converse witness of the closure
note ($M=\{e,a,b\}$ with $a\cdot a=b$, $a\cdot b=b\cdot a=a$, $b\cdot b=b$ --- i.e.\ $M=\{e,a,a^2\}$
with $a^3=a$ --- and $K(e)=e,\,K(a)=a,\,K(b)=a$)
has $\simK$-classes $\{e\},\{a,b\}$, a two-sided congruence, so \mLA{} holds; yet
$\Img K=\{e,a\}$ is not closed because $a\cdot a=b\notin\Img K$. So \emph{\mLA{} does not imply
closure}. The two conditions are incomparable.
\end{proof}

The failure in the first separator is diagnostic. Closure says the \emph{image} is a subcategory;
\mLA{} says the \emph{kernel} $\simK$ is a left congruence. Here the classes are $\{0,1\}$ and $\{2\}$,
and left-multiplying the class $\{0,1\}$ by $2$ sends $2\cdot0=2$ into class $\{2\}$ but $2\cdot1=1$
into class $\{0,1\}$: left multiplication by $2$ does not respect $\simK$. A closed image is blind to
this; \mLA{} is not.

\begin{remark}[the diamond]
\label{rem:diamond}
Theorem~\ref{thm:strict} draws a diamond over associativity:
\[
\begin{array}{ccc}
 & \Norm_K\text{ associative }=\RA{} & \\[2pt]
 \nearrow & & \nwarrow \\[2pt]
 \mLA{}:\ \simK\text{ a left congruence} & & \text{closure: }\Img K\text{ a subcategory}\\[2pt]
 \text{[\mLA{} holds, not closed]} & & \text{[closed, \mLA{} fails]}
\end{array}
\]
Each upward arrow is strict and the two upper conditions are logically independent. Over groups
(§\ref{sec:thm3}) they are the two classical mechanisms: $\mLA{}\Leftrightarrow H$ normal, closure
$\Leftrightarrow T$ a subgroup. Handedness $=\LA{}\wedge\neg\RA{}$ therefore requires both $\neg\mLA{}$
and $\neg$closure --- each is necessary for associativity to fail, neither alone is the obstruction.

One caution: the bottom vertex is \emph{not} the join of the two upper ones. Associativity (\RA{}) is
strictly weaker than ``$\mLA{}$ or closure'': there are normalizers that are associative with
\emph{neither} \mLA{} nor closed image. The smallest such lie at order~$4$ among monoids (six
normalizers, found by exhaustive search; §\ref{sec:verify}) and at $S_3\times C_2$ among groups. Only
at order~$\le3$ do the two coincide --- there associativity $=\mLA{}\vee\text{closure}$ exactly, with
zero exceptions --- which is why the two-condition picture looks complete on the small examples but is
not the general story. The general account is Proposition~\ref{prop:tfae}.
\end{remark}

\section{Over a group: the mirror law is normality}
\label{sec:thm3}

Let $G$ be a group, $H\le G$ a subgroup, and $T$ a \emph{right transversal} (one representative per
right coset $Hg$, with $\id\in T$). Let $K\colon G\to G$ send $g$ to the unique $t\in T$ with $Hg=Ht$;
then $\Img K=T$ and $\sh$ is the transversal (right-loop) product $t\sh t'=K(t\cdot t')$.

Two facts are immediate. First, \LA{} holds for \emph{every} transversal rep map: since $K(x)\in Hx$,
we have $H\cdot K(x)=H\cdot x$, hence $H\cdot(K(x)\cdot c)=H\cdot(x\cdot c)$ and
$K(K(x)\cdot c)=K(x\cdot c)$. Second, $K(a)=K(b)\Leftrightarrow Ha=Hb$, so the classes of $\simK$ are
exactly the \emph{right cosets} of $H$.

\begin{theorem}
\label{thm:normal}
With $K$ the right-coset-representative map of $H\le G$: $\mLA{}$ holds $\Leftrightarrow H\lhd G$.
When $H\lhd G$, Theorem~\ref{thm:quotient}(3) specializes to the classical isomorphism
$(T,\sh)\cong G/H$.
\end{theorem}

\begin{proof}
By Remark~\ref{rem:leftcong}, $\mLA{}\Leftrightarrow\simK$ is a left congruence. The $\simK$-classes
are right cosets, so left-compatibility reads: $Ha=Hb\Rightarrow H(ca)=H(cb)$ for all $c$.
Equivalently, for all $a,b,c$: $ab^{-1}\in H\Rightarrow(ca)(cb)^{-1}=c(ab^{-1})c^{-1}\in H$. Letting
$ab^{-1}$ range over all of $H$ and $c$ over $G$, this is exactly $cHc^{-1}\subseteq H$ for all $c$,
i.e.\ $H\lhd G$. When $H$ is normal, $\simK$ is equality of (two-sided) cosets, so
$\mathcal{A}/\!\simK=G/H$ and $\Phi\colon t\mapsto Ht$ is the group isomorphism $(T,\sh)\cong G/H$.
\end{proof}

So the normal-subgroup family of associative transversals \emph{is} the \mLA{} family. The remaining
associative transversals --- associative with $H$ not normal --- are governed by an explicit ternary
criterion.

\begin{theorem}[the exact boundary]
\label{thm:criterion}
For $f,g,h\in T$ write the de-normalization $g\cdot h=\varphi(g,h)\cdot(g\sh h)$ with
$\varphi(g,h):=(g\cdot h)(g\sh h)^{-1}\in H$. Then $(T,\sh)$ is associative $\Leftrightarrow
f\,\varphi(g,h)\,f^{-1}\in H$ for all $f,g,h\in T$.
\end{theorem}

\begin{proof}
Since $g\sh h=K(g\cdot h)\in Hg\cdot h$, indeed $\varphi(g,h)\in H$ and
$g\sh h=\varphi(g,h)^{-1}\cdot(g\cdot h)$. By \LA{} and Theorem~A of the positive-home note,
$(f\sh g)\sh h=K(f\cdot g\cdot h)$; and
\[
  f\sh(g\sh h)=K\big(f\cdot(g\sh h)\big)=K\big(f\cdot\varphi(g,h)^{-1}\cdot g\cdot h\big).
\]
Associativity is $K(f\cdot g\cdot h)=K(f\cdot\varphi(g,h)^{-1}\cdot g\cdot h)$ for all $f,g,h\in T$.
Using $K(x)=K(y)\Leftrightarrow xy^{-1}\in H$, this is
\[
  (f\cdot g\cdot h)\big(f\cdot\varphi(g,h)^{-1}\cdot g\cdot h\big)^{-1}
   = f\cdot g\cdot h\cdot(g\cdot h)^{-1}\cdot\varphi(g,h)\cdot f^{-1}
   = f\,\varphi(g,h)\,f^{-1}\in H
\]
for all $f,g,h\in T$.
\end{proof}

\begin{remark}[core form: the $f$-quantifier disappears]
\label{rem:core}
Because $G=HT$ and $(hf)^{-1}H(hf)=f^{-1}Hf$ for $h\in H$, running $f$ over the transversal $T$ is the
same as running it over all of $G$, so $\bigcap_{f\in T}f^{-1}Hf=\bigcap_{g\in G}g^{-1}Hg
=\mathrm{core}_G(H)$, the \emph{normal core} --- the largest normal subgroup of $G$ contained in $H$.
Thus ``$f\varphi(g,h)f^{-1}\in H$ for all $f$'' is just $\varphi(g,h)\in\mathrm{core}_G(H)$, and
Theorem~\ref{thm:criterion} becomes
\[
  (T,\sh)\text{ associative}\ \Longleftrightarrow\ \varphi(T\times T)\subseteq\mathrm{core}_G(H)
  \ \Longleftrightarrow\ \text{the image of }T\text{ in }G/\mathrm{core}_G(H)\text{ is a subgroup.}
\]
The two ends of the diamond are the two extremes of the core: \mLA{} is $\mathrm{core}_G(H)=H$ (i.e.\
$H\lhd G$), and closure is $\varphi\equiv\id$, where the core plays no role. This reads off the
enumeration. In $S_3$ every proper subgroup is either normal ($A_3$) or has \emph{trivial} core (the
order-$2$ subgroups), and trivial core forces $\varphi\equiv\id$ on any associative transversal
--- i.e.\ closure --- so $S_3$ has \emph{no} strictly-between transversal; the first groups with a
strictly-between regime are those with a non-trivial proper core, e.g.\ $S_3\times C_2$ (core $C_2$),
which has nine.
\end{remark}

\begin{remark}[the criterion is the diamond]
\label{rem:unify}
The criterion $f\varphi(g,h)f^{-1}\in H$ is satisfied in two canonical ways, matching
Remark~\ref{rem:diamond}:
\begin{itemize}
  \item \emph{closure / $T$ a subgroup:} then $g\cdot h\in T$, so $g\sh h=g\cdot h$ and
        $\varphi(g,h)=\id$; the criterion is vacuous.
  \item \emph{normality / \mLA{}:} $\varphi(g,h)\in H$ and $H\lhd G$, so $f\varphi(g,h)f^{-1}\in H$
        automatically.
\end{itemize}
The criterion is strictly weaker than either: it permits a nontrivial $\varphi$, valued in a
non-normal $H$, that merely stays in the normal core $\mathrm{core}_G(H)$ (Remark~\ref{rem:core}).
That strictly-between regime is the honest content of ``associative without normality or closure''.
\end{remark}

\section{The honest unifier: one condition, two specializations}
\label{sec:unifier}

Sections~\ref{sec:thm2}--\ref{sec:thm3} might read as though, over a \emph{general} associative
category $\mathcal{A}$, the two sufficient conditions --- \mLA{} and closure of $\Img K$ --- were the
two faces of one mechanism. They are not. Theorem~\ref{thm:strict} shows they are incomparable, and
off the group setting there is no $\varphi$ to carry a common criterion: ``one mechanism and its two
classical degenerations'' is a group phenomenon (§\ref{sec:thm3}), not a general theorem. What
\emph{is} general is the following, isolated by the referee (uid~239) and adopted here as the honest
statement of ``one condition behind both''. It pins down the exact restriction of \mLA{} that is
equivalent to associativity.

Write $\theta_K$ for the \emph{largest two-sided congruence contained in $\simK$}; explicitly,
\[
  a\mathbin{\theta_K}b\ :\Longleftrightarrow\ K(x\cdot a\cdot y)=K(x\cdot b\cdot y)
  \quad\text{for all composable }x,y.
\]
(Taking $x=y=\id$ shows $\theta_K\subseteq\simK$; the two quantifiers make $\theta_K$ a two-sided
congruence; and any two-sided congruence inside $\simK$ is contained in $\theta_K$, so it is the
largest.)

\begin{proposition}[the honest unifier]
\label{prop:tfae}
Let $K$ be a normalizer (so \LA{} holds). The following are equivalent.
\begin{enumerate}
  \item $\Norm_K$ is associative (equivalently, \RA{} holds).
  \item \mLA{} holds \emph{on products of normalized arrows:} $K\big(c\cdot K(x)\big)=K(c\cdot x)$ for
        all $c\in\mathcal{A}$ and all $x\in\Img K\cdot\Img K$ (i.e.\ $x=g\cdot h$, $g,h\in\Img K$).
  \item $K(g\cdot h)\mathbin{\theta_K}(g\cdot h)$ for all $g,h\in\Img K$.
\end{enumerate}
\end{proposition}

\begin{proof}
\emph{(2)$\Rightarrow$(1).} Take $c=f\in\Img K$ and $x=g\cdot h$ with $g,h\in\Img K$; then (2) reads
$K\big(f\cdot K(g\cdot h)\big)=K(f\cdot g\cdot h)$, which is \RA{}. So $\Norm_K$ is associative by
Proposition~\ref{prop:thmB}.

\emph{(1)$\Rightarrow$(2).} For $c\in\mathcal{A}$ and $g,h\in\Img K$,
\[
  K\big(c\cdot K(g\cdot h)\big)
  \overset{\LA{}}{=}K\big(K(c)\cdot K(g\cdot h)\big)
  \overset{\RA{}}{=}K\big(K(c)\cdot g\cdot h\big)
  \overset{\LA{}}{=}K(c\cdot g\cdot h),
\]
the outer steps being \LA{} (read backwards and forwards) and the middle step \RA{} applied to
$K(c),g,h\in\Img K$, available because (1)$\Leftrightarrow$\RA{}.

\emph{(2)$\Leftrightarrow$(3).} Since $K$ is idempotent, $K(g\cdot h)\simK(g\cdot h)$ always; and every
normalizer has the \LA{} (right-congruence) half, i.e.\ $K(a\cdot y)=K(K(a)\cdot y)$. Condition~(3)
unfolds to $K\big(x\cdot K(g\cdot h)\cdot y\big)=K\big(x\cdot(g\cdot h)\cdot y\big)$ for all $x,y$.
Setting $y=\id$ gives (2) with $c=x$. Conversely, given (2), for all $x,y$,
\[
  K\big(x\cdot K(g\cdot h)\cdot y\big)=K\big(K(x\cdot K(g\cdot h))\cdot y\big)
  =K\big(K(x\cdot g\cdot h)\cdot y\big)=K\big(x\cdot(g\cdot h)\cdot y\big),
\]
using right-compatibility at the outer equalities and (2) (with $c=x$) at the middle. So (3) holds.
\end{proof}

\begin{remark}[what (2) is, and what it is not]
\label{rem:tfae-ends}
Condition~(2) is the honest ``single condition behind both'': it is \RA{} rewritten as \mLA{}
restricted to products of normalized arrows. The full mirror law \mLA{} (all $x\in\mathcal{A}$) and
closure of $\Img K$ are its two \emph{sufficient specializations} --- full \mLA{} gives (2) a
fortiori; closure makes (2) \emph{vacuous} (then $x=g\cdot h\in\Img K$, so $K(x)=x$ and both sides are
$K(c\cdot x)$). They are \emph{not} two degenerations of one criterion over general $\mathcal{A}$
(Theorem~\ref{thm:strict}: incomparable), but two sufficient routes to the \emph{same}
necessary-and-sufficient condition~(2)/(3). Over a group with $K$ the coset-representative map,
$\theta_K$ is the coset relation of $\mathrm{core}_G(H)$ and (3) is Theorem~\ref{thm:criterion} in the
core form of Remark~\ref{rem:core}. The equivalence (1)$\Leftrightarrow$(3) was checked exhaustively
on all $1578$ normalizers of all $156$ order-$4$ monoids (unit fixed) and all monoids of order~$\le3$,
with zero mismatches.
\end{remark}

\section{Verification and scope}
\label{sec:verify}

\paragraph{Lean (finite core, axiom-free).} The leaf \texttt{lean-mirror-LA-oracle}
(\texttt{lean/\allowbreak2026-10-22-\allowbreak mirror-LA-\allowbreak oracle.lean}, Lean~4,
Mathlib-free) machine-checks the decidable
content: on the order-$4$ handed model, \mLA{} fails (as it must, being handed); \mLA{} forces \RA{}
and $\sh$-associativity over \emph{all} monoids of order $\le3$ and all idempotent unit-fixing
$K$\footnote{This Lean range includes idempotent unit-fixing $K$ that do \emph{not} satisfy \LA{},
which the hand proof of Theorem~\ref{thm:quotient}(1) does not cover (its passage \RA{}$\Rightarrow$
associative uses \LA{}, via Proposition~\ref{prop:thmB}). The certified implication
``\mLA{}$\Rightarrow\sh$-associative'' is nonetheless true on the whole range: at order~$3$ there are
exactly $6$ such $K$ with \mLA{} but without \LA{}, and all $6$ are associative. For those the
statement is a verified computational fact, not a consequence of the hand argument.}
(zero counterexamples); $\mLA{}\Leftrightarrow\simK$ a left congruence over the same range (zero
mismatches, Remark~\ref{rem:leftcong}); the strictness separator exists at order~$3$ and not below;
and the explicit separator of §\ref{sec:thm2} (closed image, \RA{} true, \mLA{} false at $c=2,x=1$)
is certified. The bundled theorem reports \emph{no} axiom dependence --- not even \texttt{propext};
every \texttt{decide} reduces to \texttt{rfl}. We cite it rather than reproduce it.

\paragraph{Exhaustive enumeration.} \texttt{scratch/2026-10-22-mirror-LA-verify.py} (exact integer /
permutation arithmetic). Over all monoids of order $\le3$ and all their normalizers:
$\mLA{}\Rightarrow\RA{}$, $\mLA{}\Rightarrow$ associative, and $\mLA{}\Leftrightarrow\simK$
left-congruence hold with zero exceptions; right-congruence never fails under \LA{}; and, with the
unit fixed ($11$ labelled order-$3$ tables), exactly $2$ separators $\RA{}\wedge\neg\mLA{}$ occur,
both at order~$3$ (none at order $\le2$); allowing any element to be the unit ($33$ tables) the count
is $6$. Over all $(H,T)$
transversal pairs of $S_3$ ($15$) and $S_4$ ($35\,416$): \LA{} never fails,
$[\mLA{}\Leftrightarrow H\text{ normal}]$ has zero mismatches, and
$[(T,\sh)\text{ associative}\Leftrightarrow\text{the criterion}]$ has zero mismatches. The explicit
$S_3$ witness of the ``closed but \mLA{} fails'' corner is $H=\langle(1\,2)\rangle$ (order~$2$,
non-normal), $T=A_3$ (a subgroup, hence closed): \LA{}~\checkmark, \mLA{}~fails, $(T,\sh)\cong C_3$
associative.

\paragraph{Honest scope.} Theorems~\ref{thm:quotient}--\ref{thm:criterion} are proved by hand
(elementary: \textsf{(K1)}, \LA{}, \mLA{}, associativity of $\cdot$, and, for §\ref{sec:thm3}, the
right-coset structure); the enumerations corroborate but are not the proof. Minimality of the order-$3$
separator is a computational fact (exhaustive over order $\le3$), as in the closure note, and nothing
downstream depends on it. The group reading assumes a right transversal with $\id\in T$; the
left-transversal reading swaps $\LA{}\leftrightarrow\mLA{}$ and gives the mirror statements. The group
specialization (Theorem~\ref{thm:normal}) and the $S_3/S_4$ enumerations are Python-only (permutation
searches, not reducible kernel \texttt{decide}s at that scale); the Lean leaf certifies the abstract
finite core.

\section*{What this settles}

There is one honest condition behind the two associative families, and it is \emph{not} a single
mechanism with two degenerations --- over a general $\mathcal{A}$, \mLA{} and closure are incomparable
(Theorem~\ref{thm:strict}). The common condition is associativity itself, read as the restricted
mirror law of Proposition~\ref{prop:tfae}: $K(c\cdot K(x))=K(c\cdot x)$ for $x$ a product of
normalized arrows, equivalently $K(g\cdot h)\mathbin{\theta_K}(g\cdot h)$. Full \mLA{} and closure of
$\Img K$ are its two \emph{sufficient specializations} --- the first makes the kernel a two-sided
congruence, so $\Norm_K\cong\mathcal{A}/\!\simK$ is the honest quotient; the second makes the
restricted law vacuous. \mLA{} is itself a legible, checkable sufficient condition (``$\simK$ is a
left congruence''), strictly stronger than bare \RA{} and incomparable with closure; the two sit in a
diamond over associativity whose bottom vertex is strictly below their join. Over groups, \mLA{} is
normality and closure is subgroup-hood, and both solve the single conjugation criterion
$f\varphi(g,h)f^{-1}\in H$ --- equivalently $\varphi(T\times T)\subseteq\mathrm{core}_G(H)$
(Remark~\ref{rem:core}) --- which exhibits, between them, the genuine associative-without-either
regime. This is the content SHOULD~S1 asked for.

\begin{thebibliography}{9}

\bibitem{positivehome}
MacBeth, \emph{The positive home of the normalization defect: a one-sided-absorbing idempotent
normalizer} (program note, 2026-10-19; registry node
\texttt{positive-home-one-sided-absorbing-idempotent}). Source of Theorems~A and~B and the handedness
definition recalled here.

\bibitem{triangulation}
MacBeth, \emph{What normalization's non-associativity is, and is not: three obstruction theories that
miss it, and the idempotent normalizer that receives it} (triangulation unit,
\texttt{expository/2026-10-21-normalization-defect-triangulation}, content commit \texttt{6df0043}).
Contains the closure criterion and the provenance discussion (Baer; Lal; Kiechle; Grabowski).

\bibitem{dlrs}
E.~Di~Lavore, M.~Rom\'an, P.~Sz\'ekely, \emph{Normalization is a sesquilaw}, arXiv:2510.01131. Origin
of the magmoid $\Norm$ and the renormalization idempotent.

\bibitem{baer}
R.~Baer, Nets and groups, \emph{Trans.\ Amer.\ Math.\ Soc.} \textbf{46} (1939) 110--141.

\bibitem{lal}
R.~Lal, Transversals in groups, \emph{J.\ Algebra} \textbf{181} (1996) 70--81.

\bibitem{kiechle}
H.~Kiechle, \emph{Theory of K-Loops}, Springer LNM~1778 (2002).

\end{thebibliography}

\noindent\footnotesize The oidified (categorical) version of this transversal/right-loop construction
is published (Grabowski, \emph{An introduction to loopoids}): see the provenance discussion of
the triangulation unit~[2], whose source index logs it. It is not re-cited here directly pending a
deep-read registration.

\end{document}
